\section{Augustine and Metaphysical Positivism}

\begin{quotex}
Noli foras ire, in te ipsum redi: in interiore hominis habitat veritas.
\flright{\textsc{St. Augustine}}
\end{quotex}


In \textbf{Auguste Comte}'s explanation of Positivism, only data that could be experienced by the senses was acceptable. From that, laws would be inferred, free from any metaphysical or ideological presuppositions. That is sound practice for empirical science, but not for spiritual development. \textbf{Julius Evola} proposed a metaphysical positivism\footnote{\url{https://gornahoor.net/?p=1189}} which would be based, not on philosophical speculations, but on direct, intuitive spiritual experience.

We continue to receive hate mail complaining that we criticize someone's favourite cause unfairly. These e-mails are simply confusions built on incomprehension. If only it were a matter as easy as putting the on the right team jersey, or adopting some man-made thought scheme and moral code. Neither is it a matter of being born with the ``right'' genetic structure, as though no personal efforts were required.

No, it is not a matter of playing the game, if only for the correct team. Nor is it about sitting out the game and not playing. Rather the task is to transcend the game, while observing it, making sense of the hidden forces guiding the movements of the players, and participating on a super-conscious level. We have several times mentioned various Hermetic meditations and trials. Here are a few to try this week.


\begin{description}

\item[Empty Mind]

For this exercise, purge the mind of every theory, creed, belief, custom, habit, or old wives tales. Then eliminate any desires, fantasies, or other emotional disturbance. What is left over? What rushes in to fill the void? What if you were with others in a similar state? What kind of world would you create?
\item[Amor Fati]

Think about your life as it is and how it came to be that way. Forget about those who, you believe, frustrated or disappointed you. Can you love your life such as it is? Would you be willing to live it over again with everything happening in the exact same way? Can you accept ownership for every aspect in your life? Are you aware of any transcendent forces guiding your life?
\item[Cemetery]

Make a visit to a cemetery, preferably one that has significance to you. Pay your respects and be quiet. You may not learn about the dead, but you may learn something about yourself.
\end{description}


\paragraph{Augustine of Hippo}


We have claimed that \textbf{Augustine of Hippo} (along with Plotinus and Boethius) was one of the primary intellectual forces behind what Evola calls the Nordic-Roman Middle Ages and we have been calling Christendom. The conventional historico-theological approach is to analyze his theological positions vis-à-vis the various issues of his time, such as grace, predestination, salvation, and so on. But this is only at the level of thought and is impossibly distant from a ``metaphysical positivism''.

The Orthodox priest, Matthew Rafael Johnson, has made available a lecture on Augustine\footnote{\url{https://gornahoor.net/library/StAugustine.mp3}} that is more in line with a metaphysical positivism (we do not imply that he would agree with much on Gornahoor nor vice versa). Augustine wrote: ``Don't go outside yourself, turn back to yourself. Truth resides in the depths of man.'' That is always the starting point. Dogma can arise only after the fact of that direct inner experience of Truth. Fr. Johnson explains the inner states that Augustine is striving to convey in his writings. An interesting point is that Fr. Johnson returns to the the idea of St Anthony of the Desert\footnote{\url{https://gornahoor.net/?p=501}} who sees an insurmountable ontological divide between men born again of the spirit and those without spirit. The latter have no eternal life, perhaps a more more shadowy existence known to the ancients.\footnote{\url{https://gornahoor.net/?p=2507\#more-2507}}


\flrightit{Posted on 2011-08-04 by Cologero}

\begin{center}* * *\end{center}

\begin{footnotesize}
\begin{sffamily}
\texttt{Matt on 2011-08-04 at 23:20 said:}

Funny enough, I listened to that lecture this past saturday. Fr. Johnson is the only person on that site I listen to.

\hfill

\texttt{Cologero on 2011-08-05 at 00:04 said:}

Good judgment about that site, Matt. Unlike Fr. Johnson, Orthodox writers show little appreciation for Augustine, although Seraphim Rose wrote a lukewarm book about him. Johnson seems to me to be very nearly a Marcionite, barely bringing in the OT and almost replacing it with Platonism.

\hfill

\texttt{Auld Wat on 2015-12-19 at 09:48 said:}

There is an error with the audiofile of the lecture!

\hfill

\texttt{Cologero on 2015-12-19 at 11:31 said:}

Auld Wat, try this link:
\url{http://www.gornahoor.net/library/StAugustine.mp3}.


\hfill

\texttt{Kevin on 2017-08-05 at 14:37 said:}

Absolute excellence!!! As a student, I am in continual need of reminding and exhortation --- habit is difficult to change.


\hfill

\end{sffamily}
\end{footnotesize}

